\documentclass[11pt,a4paper]{amsart}

\usepackage[T1]{fontenc}
\usepackage{lmodern}
\usepackage[protrusion=true,expansion=false]{microtype}
\usepackage[margin=30mm]{geometry}
\usepackage{amsmath,amssymb,amsthm,mathtools}
\usepackage{enumitem}
\usepackage[dvipsnames]{xcolor}
\usepackage[colorlinks=true,linkcolor=MidnightBlue,citecolor=MidnightBlue,urlcolor=MidnightBlue]{hyperref}
\hypersetup{
  pdftitle={Logarithmic Residual-Complexity Barriers for Local Mobius Cancellation},
  pdfauthor={Mayk Loide Baccaro},
  pdfsubject={Necessary complexity and pair-compression theorems for local Mobius cancellation schemes}
}

\newtheorem{theorem}{Theorem}[section]
\newtheorem{proposition}[theorem]{Proposition}
\newtheorem{corollary}[theorem]{Corollary}
\newtheorem{lemma}[theorem]{Lemma}
\theoremstyle{definition}
\newtheorem{definition}[theorem]{Definition}
\theoremstyle{remark}
\newtheorem{remark}[theorem]{Remark}

\newcommand{\N}{\mathbb N}
\newcommand{\R}{\mathbb R}
\newcommand{\mob}{\mu}
\newcommand{\om}{\omega}
\newcommand{\M}{M}
\setlength{\emergencystretch}{2em}

\title[Residual complexity for local M\"obius cancellation]
{Logarithmic Residual-Complexity Barriers\\
for Local M\"obius Cancellation}

\author[Mayk Loide Baccaro]{Mayk Loide Baccaro\\
Independent researcher\\
\href{mailto:m.baccaro7663@student.nu.edu}{\texttt{m.baccaro7663@student.nu.edu}}}

\date{18 August 2026}
\subjclass[2020]{Primary 11N37; Secondary 05C70, 11A25}
\keywords{M\"obius function, squarefree integers, matching, residual complexity, Mertens function, hypergraph compression}

\begin{document}
\raggedbottom

\begin{abstract}
Consider cancellation certificates for the M\"obius function that pair
squarefree integers after removing their greatest common divisor.  If the
residual cofactors have at most $R(T)$ prime factors at scale
$T=\max(a,b)$, every edge has length $O(T^\alpha)$, and only
$O(X^\beta)$ vertices through $X$ remain unmatched, then
\[
 \limsup_{X\to\infty}\frac{R(X)\log\log X}{\log X}
 \geq (1-\alpha)(1-\beta).
\]
The square-root scale therefore forces a limsup of at least $1/4$.
Finite vertex-disjoint, unit-weight, zero-sum M\"obius cells with a common
core compress exactly to opposite-sign pairs.  When the residual budget
depends on scale, the compression loses only endpoints near jumps of $R$.
Thus short local certificates require logarithmically growing residual
complexity, and common-core cells obey the same pairwise limitation.  The
results construct no matching and give no estimate for the actual Mertens
function.
\end{abstract}

\maketitle
\enlargethispage{12pt}

\section{Introduction}

Let $\mob$ be the M\"obius function and
\[
 \M(X)=\sum_{n\leq X}\mob(n).
\]
The Riemann hypothesis is equivalent to
$\M(X)=O_\varepsilon(X^{1/2+\varepsilon})$ for every $\varepsilon>0$;
see Titchmarsh~\cite[Chapter XIV]{Titchmarsh1986} or
Montgomery--Vaughan~\cite[Chapter 15]{MontgomeryVaughan2007}.  Since
$\mob(n)\in\{-1,0,1\}$, a combinatorial route would pair most squarefree
integers of opposite M\"obius sign with short edges and charge the summatory
error to unmatched vertices and edges crossing the cutoff.

We study the necessary nonlocality of a broad gcd-residual matching and a
natural enlargement from pairs to finite cells.  Related work treats sign
patterns of $\mob$ and the Liouville function by analytic methods
\cite{MatomakiRadziwillTao2016}, together with harmonic and matrix structures
on the squarefree integers~\cite{Hilberdink2014}.  McNew and Pomerance study
coprime matchings between the divisor set of an integer and an initial
interval, proving in particular that every squarefree integer is matchable
\cite{McNewPomerance2026}.  Their edge condition concerns coprimality.  The
edges here must have opposite M\"obius signs, short displacement, and a bounded
residual after the common divisor is removed.  The controlling question is
the obstruction forced by the retained greatest common divisor in such an
edge.

Aymone and Liu study analytic bounds for summatory functions of
multiplicative functions resembling M\"obius, including a separate
one-quarter lower-bound exponent in Liu's class~\cite{Aymone2020,Liu2023}.
Their objects are altered multiplicative functions; the objects here are
local matching certificates for the M\"obius function itself.

Products of many primes in a short dyadic band form a large same-sign family.
A small residual budget forces any repair
edge incident to that family to retain a greatest common divisor larger than
the permitted displacement.  This gives the product threshold
$(1-\alpha)(1-\beta)$ in Theorem~\ref{thm:threshold}.

The second part tests a larger zero-sum cell.  Distinct squarefree vertices
with unit M\"obius weights and one common greatest common divisor can be
rank-paired by sign.  The pairing preserves the exact signed boundary mass at
every cutoff.  If the budget changes with the endpoint, failed derived pairs
are confined to neighborhoods of its jump set; Theorem~\ref{thm:sparse-jump}
quantifies the loss.

The scope is the stated certificate class.  Weighted flows, overlapping
cells, non-common-core structures, cross terms, signed residual estimates,
and dense-jump budgets require separate arguments.  The conclusions give no
lower bound for $|\M(X)|$ and do not cover every combinatorial approach to
M\"obius cancellation.

\section{The maximal residual pair class}

Write $\om(n)$ for the number of distinct prime divisors of $n$.  Let
$R:[1,\infty)\to\N\cup\{0\}$ be nondecreasing.

\begin{definition}[Residual-admissible pair]\label{def:pair}
Distinct squarefree integers $a,b$ form an $R$-admissible pair if, with
\[
 g=(a,b),\qquad T=\max(a,b),
\]
one has
\begin{equation}\label{eq:budget}
 \om(a/g)\leq R(T),\qquad \om(b/g)\leq R(T),
\end{equation}
and $\om(a/g)$ and $\om(b/g)$ have opposite parity.
\end{definition}

The parity condition gives $\mob(a)=-\mob(b)$, since $g,a/g,b/g$ are
squarefree.  Every such pair also obeys
\begin{equation}\label{eq:gcd-floor}
 |a-b|=g\,|a/g-b/g|\geq g.
\end{equation}
Definition~\ref{def:pair} permits every edge satisfying the stated budget,
so any extra scheduling or graph restriction can only remove edges.

\begin{lemma}[Matching identity]\label{lem:matching}
Let $\mathcal P$ be a finite family of disjoint opposite-sign pairs and let
$U(X)$ be the number of squarefree integers at most $X$ not covered by
$\mathcal P$.  If $C(X)$ is the number of pairs with exactly one endpoint at
most $X$, then
\[
 |\M(X)|\leq U(X)+C(X).
\]
\end{lemma}

\begin{proof}
Pairs lying completely below or completely above $X$ contribute zero.  Each
crossing pair and each uncovered squarefree vertex contributes a number of
absolute value one.  Integers with $\mob(n)=0$ contribute nothing.
\end{proof}

\section{The logarithmic complexity threshold}

\begin{theorem}[Residual-complexity barrier]\label{thm:threshold}
Fix $0\leq\alpha<1$, $0\leq\beta<1$, and $C>0$.  Suppose a matching of
squarefree integers uses only $R$-admissible pairs and satisfies
\begin{equation}\label{eq:short}
 |a-b|\leq C\max(a,b)^\alpha
\end{equation}
on every used edge.  If the number of unmatched squarefree vertices through
$X$ is $O(X^\beta)$, then
\begin{equation}\label{eq:threshold}
 \limsup_{X\to\infty}
 \frac{R(X)\log\log X}{\log X}
 \geq (1-\alpha)(1-\beta).
\end{equation}
\end{theorem}

\begin{proof}
Assume that the limsup in~\eqref{eq:threshold} is
\[
 r<(1-\alpha)(1-\beta).
\]
Choose
\begin{equation}\label{eq:delta}
 \frac{r}{1-\alpha}<\delta<1-\beta.
\end{equation}
For large $y$, put $k=\lfloor y^\delta\rfloor$, let
$m_y=\pi(2y)-\pi(y)$, let $S_y$ contain the products of $k$ distinct primes
in $(y,2y]$, and put $X_y=(2y)^k$.

The prime number theorem in a fixed dyadic interval
\cite{MontgomeryVaughan2007} gives
$m_y\gg y/\log y$.  Since $k=y^{\delta+o(1)}=o(m_y)$, the standard upper and
lower binomial estimates give
\begin{equation}\label{eq:family}
 |S_y|=\binom{m_y}{k}=X_y^{1-\delta+o(1)}.
\end{equation}
Every element of $S_y$ has M\"obius sign $(-1)^k$.  The choice of $\delta$
in~\eqref{eq:delta} makes this family larger than $X_y^\beta$ by a fixed
power for all sufficiently large $y$.

Fix $n\in S_y$ and suppose that it has a partner $b$.  If $b\geq2n$, then
\eqref{eq:short} gives $b^{1-\alpha}\leq2C$, which fails for large $y$.
Thus $b<2n\leq2X_y$, and
\begin{equation}\label{eq:logscale}
 \frac{\log(2X_y)}{\log\log(2X_y)}
 =\left(\frac1\delta+o(1)\right)k.
\end{equation}
The assumed limsup bound, monotonicity of $R$, and~\eqref{eq:delta} yield
some $\eta>0$ such that
\begin{equation}\label{eq:R-bound}
 R(\max(n,b))\leq R(2X_y)<(1-\alpha-\eta)k
\end{equation}
for all sufficiently large $y$.

Let $g=(n,b)$.  The residual $n/g$ contains at most $R(2X_y)$ of the $k$
prime factors of $n$.  Combining the retained prime factors with
\eqref{eq:gcd-floor} and the short-edge bound gives
\begin{equation}\label{eq:key-sandwich}
 y^{(\alpha+\eta)k}
 < y^{k-R(2X_y)}
 \leq g
 \leq |n-b|
 \leq C(2X_y)^\alpha
 =y^{(\alpha+o(1))k},
\end{equation}
which is impossible for large $y$.  Hence every member of $S_y$ is
unmatched.  Equation~\eqref{eq:family} contradicts the asserted
$O(X_y^\beta)$ unmatched count.
\end{proof}

\begin{corollary}[Square-root matching scale]\label{cor:quarter}
Fix $R$.  Suppose that for every $\varepsilon>0$ there is a
residual-admissible matching, allowed to depend on $\varepsilon$, whose edge
lengths and unmatched counts are
\[
 O_\varepsilon(X^{1/2+\varepsilon}).
\]
Then necessarily
\[
 \limsup_{X\to\infty}
 \frac{R(X)\log\log X}{\log X}\geq\frac14.
\]
\end{corollary}

\begin{proof}
Apply Theorem~\ref{thm:threshold} with
$\alpha=\beta=1/2+\varepsilon$ and let $\varepsilon\downarrow0$.
\end{proof}

A matching construction additionally requires Hall-type expansion or another
mechanism for controlling collisions.

\section{Common-core cells collapse to pairs}

\begin{definition}[Common-core M\"obius cell]
A cell is a finite set $H$ of distinct squarefree positive integers with
\[
 \sum_{n\in H}\mob(n)=0.
\]
Write $g_H=\gcd\{n:n\in H\}$, $T_H=\max H$, and
\[
 S_H(X)=\sum_{\substack{n\in H\\ n\leq X}}\mob(n).
\]
\end{definition}

\begin{theorem}[Exact fixed-budget pair compression]\label{thm:compression}
Suppose a common-core M\"obius cell $H$ satisfies
\[
 \om(n/g_H)\leq R_H\qquad(n\in H).
\]
List its positive vertices increasingly as $p_1,\ldots,p_m$ and its negative
vertices increasingly as $q_1,\ldots,q_m$.  Pair $p_i$ with $q_i$.  Then
every derived pair $\{a,b\}$ satisfies
\[
 \mob(a)=-\mob(b),\quad
 \om(a/(a,b)),\om(b/(a,b))\leq R_H,\quad
 |a-b|\leq\operatorname{diam}(H).
\]
For every real $X$, exactly $|S_H(X)|$ of these pairs cross $X$.
\end{theorem}

\begin{proof}
The zero-sum hypothesis gives the same number of positive and negative
vertices.  For a paired $a,b$, the common core $g_H$ divides $h=(a,b)$.
Since the integers are squarefree,
\[
 \om(a/h)\leq\om(a/g_H)\leq R_H,
 \qquad
 \om(b/h)\leq\om(b/g_H)\leq R_H.
\]
Opposite M\"obius signs give opposite residual parities.  The diameter bound
is immediate.

At a cutoff $X$, let $P(X)$ and $Q(X)$ count the positive and negative
endpoints through $X$.  Increasing rank pairing matches the common initial
portion.  The crossing-pair count is therefore
$|P(X)-Q(X)|=|S_H(X)|$.
\end{proof}

The final identity preserves the exact signed prefix mass simultaneously at
every cutoff.

\section{Endpoint budgets and sparse jumps}

If a cell is admitted at budget $R(T_H)$, a derived pair with maximum
endpoint $m<T_H$ may exceed the smaller budget $R(m)$.  The following example
shows the endpoint issue exactly.

\begin{proposition}[Sharp finite endpoint gap]\label{prop:gap}
Let
\[
 H=\{2,5,21,33\},\qquad
 R(t)=\begin{cases}1,&t<30,\\2,&t\geq30.\end{cases}
\]
Then $H$ is a zero-sum common-core cell admitted at scale $33$, but neither
opposite-sign pair partition makes every pair admissible at its own maximum
endpoint.
\end{proposition}

\begin{proof}
The signs are $(-1,-1,+1,+1)$ and every vertex has at most two prime
factors.  There are two opposite-sign pair partitions.  In each partition,
the pair not containing $33$ has maximum below $30$ and contains an endpoint
with two residual prime factors, so it violates budget one.
\end{proof}

Let $J(Y)$ be the number of integer jump locations $j\leq Y$ with
$R(j)>R(j-1)$.

\begin{theorem}[Sparse-jump repair]\label{thm:sparse-jump}
Let $R:\N\to\N\cup\{0\}$ and $D:\N\to[0,\infty)$ be nondecreasing.  At a
cutoff $X$, let $\mathcal F_X$ be a finite vertex-disjoint family of
common-core M\"obius cells satisfying
\[
 \om(n/g_H)\leq R(T_H),\qquad
 \operatorname{diam}(H)\leq D(T_H).
\]
Rank-pair every cell as in Theorem~\ref{thm:compression}, and delete a derived
pair when it violates the residual budget at its own maximum endpoint.  If
$D(t)\leq t/2$ for all sufficiently large $t$, then, up to a fixed initial
constant $B_0$, the number $B(X)$ of deleted endpoints at most $X$ obeys
\begin{equation}\label{eq:jump-bound}
 B(X)\leq B_0+\bigl(2\lfloor D(2X)\rfloor+1\bigr)J(2X).
\end{equation}
If also $D(t)\leq Ct^\alpha$ with $\alpha<1$, every retained pair satisfies
\begin{equation}\label{eq:inherited-gap}
 |a-b|\leq C2^\alpha\max(a,b)^\alpha
\end{equation}
for all sufficiently large endpoints.
\end{theorem}

\begin{proof}
Let a deleted pair have maximum endpoint $m$.  Common-core inheritance still
bounds both residuals by $R(T_H)$; deletion therefore implies
$R(m)<R(T_H)$.  Hence there is a jump $j\in(m,T_H]$.  Every endpoint of the
pair lies within $D(T_H)$ of $j$.

If a deleted endpoint $n\leq X$ lies in $H$, then
\[
 T_H-n\leq D(T_H)\leq T_H/2,
\]
so $T_H\leq2X$ and $D(T_H)\leq D(2X)$.  Each endpoint is therefore an
integer in a radius-$D(2X)$ neighborhood of a jump through $2X$.
Vertex-disjointness ensures that no integer endpoint is counted twice, which
proves~\eqref{eq:jump-bound}.

For the retained-pair gap, put $m=\max(a,b)$.  Eventually
$m\geq T_H-D(T_H)\geq T_H/2$, and therefore
\[
 |a-b|\leq D(T_H)\leq CT_H^\alpha\leq C2^\alpha m^\alpha.
\]
\end{proof}

\begin{corollary}[Compression at square-root scale]\label{cor:compression}
Suppose $R(Y)=O(\log Y/\log\log Y)$ and, for every $\delta>0$ and
sufficiently large $X$, there is a family $\mathcal F_{\delta,X}$ as in
Theorem~\ref{thm:sparse-jump} with
\[
 D(t)=O_\delta(t^{1/2+\delta})
\]
and total uncovered plus signed cell-boundary mass
\[
 O_\delta(X^{1/2+\delta}).
\]
Here the boundary mass is the sum of $|S_H(X)|$ over the cells.  Then for
every $\varepsilon>0$, that cutoff-specific cell certificate compresses to an
opposite-sign pair certificate with unmatched plus crossing contribution
$O_\varepsilon(X^{1/2+\varepsilon})$.
\end{corollary}

\begin{proof}
An integer-valued nondecreasing $R$ has $J(Y)\leq R(Y)+O(1)$.  Choose
$0<\delta<\min(\varepsilon,1/2)$.  Equation~\eqref{eq:jump-bound} then
contributes
\[
 O_\delta\!\left(X^{1/2+\delta}
 \frac{\log X}{\log\log X}\right)
 =O_\varepsilon(X^{1/2+\varepsilon}).
\]
Theorem~\ref{thm:compression} preserves the signed boundary mass, and
\eqref{eq:inherited-gap} preserves the power exponent.
\end{proof}

The family may depend on the queried cutoff when the desired conclusion is
pointwise and all constants are uniform.  A global or online matching claim
requires one coherent family.

\section{Scope and remaining mechanisms}

Theorem~\ref{thm:threshold} rules out square-root performance for short
gcd-residual pairs with a fixed or subthreshold edit budget.
Theorems~\ref{thm:compression} and~\ref{thm:sparse-jump} transfer the same
limitation to common-core unit cells under a sparse-jump endpoint budget.

The following mechanisms remain outside the results:
\begin{enumerate}[label=(\roman*)]
\item weighted or signed flows in place of raw unit M\"obius mass;
\item overlapping cells with a proved ownership or conservation law;
\item cells without a common gcd or with aggregate residual control;
\item cross-cofactor moves and cancellation among unmatched signs;
\item dense-jump budgets with an independent complexity theorem; and
\item analytic estimates not mediated by a short local matching.
\end{enumerate}

Unmatched cardinality can exceed $|\M(X)|$ because the unmatched signs may
cancel.  The conclusions apply to the stated cancellation certificates.

\appendix
\section{Verification remarks}

The proofs are symbolic.  The threshold theorem combines the displacement
factorization~\eqref{eq:gcd-floor}, the band-product estimate
\eqref{eq:family}, the scale conversion~\eqref{eq:logscale}, and the strict
exponent separation in~\eqref{eq:key-sandwich}.  Exact pair compression
preserves $|S_H(X)|$ at every prefix.  The sparse-jump argument charges a
deleted pair only after establishing $R(m)<R(T_H)$.  Finite exact-arithmetic
checks recover the displayed cell and its two pairings.  The printed
deductions establish the asymptotic theorem.

\section*{Data and code availability}

The asymptotic arguments are self-contained and use no external dataset.  An
accompanying Lean~4 development proves the fixed-band prime supply, exact
prime-product population, pointwise matching obstruction, finite and all-scale
residual-or-unmatched dichotomy, logarithmic scale conversion, and
residual-budget separation.  It also checks finite opposite-sign cancellation,
the matching bound, zero-sum cell balance, the strict exponent contradiction,
the square-root constant, and the endpoint-doubling connector.  The final
sparse-binomial comparison giving the printed limsup theorem, exact rank
compression, and the global sparse-jump argument remain manuscript-level
proofs.

\section*{Assistance disclosure}

AI systems were used extensively in mathematical derivation, Lean proof
development, literature discovery, organization, and typesetting.  The author
remains responsible for every mathematical statement, proof, citation, and
submission decision.

\bibliographystyle{amsplain}
\bibliography{references}

\end{document}
